\section{Technical Details \textendash \ Fragmentation Index}
\label{app:FI}

In this appendix we prove Proposition~\ref{prop:fragmentation}, which relates the cost of the optimal short-term AI deployment strategy for a fixed job to the fragmentation index of that job.
Intuitively, we expect a production process in which automatable tasks are clustered together will benefit more from AI deployment than one where steps with high and low levels of automation susceptibility are interspersed.  
This is due to the potential benefits from AI chains.

First recall the definition of fragmentation index for a given step sequence $\mathcal{S} = (s_1, \dotsc, s_m)$.  
As a reminder, we assume for Proposition~\ref{prop:fragmentation} that $\AItime{i} = 1$ for all $s_i$.  
Consider a random process in which each step $s_i$ succeeds independently with probability $q_i$; any step that does not succeed is said to fail.  
Write $F$ for the set of steps that fail, and $\mathcal{C} = \{C_1, \dotsc, C_k\}$ for the random variable representing the collection of maximal connected components of non-failed steps.  
The weight of each $C_j \in \mathcal{C}$ is defined to be $\omega(C_j) = \min\{ 1, \sum_{s_i \in C_j} \manualTime{i} \}$.  
That is, each $C_j$ has weight $1$ unless the sum of the manual time costs for each of its steps is less than $1$ (due to our assumption that AI management costs are normalized to $1$).

We now introduce some notation: for each $s_i$, write $\minTime{i} = \min\left\{\manualTime{i}, \frac{\AItime{i}}{q_i}\right\}$.  That is, $\minTime{i}$ is the minimum cost to implement step $s_i$ individually (whether manually or through AI augmentation).  
Given a realization of $\mathcal{C}$ and $F$, we define the realized fragmentation to be 
\[ \sum_{s_i \in F}\minTime{i} + \sum_{C_j \in \mathcal{C}} \omega(C_j).\]
The \emph{fragmentation index} is defined to be the expected value of the realized fragmentation.

We now fix the sequence $\mathcal{S} = (s_1, \dotsc, s_m)$, write $FI$ for the fragmentation index of $\mathcal{S}$, and let $OPT$ be the cost of the optimal (i.e., time-minimizing) task structure for those steps.  

We begin by showing that $FI$ is not much larger than $OPT$.

\begin{proposition}
\label{prop:FI.upper.bound}
    $FI \leq \tfrac{5}{4}OPT$.
\end{proposition}
\begin{proof}
Fix an arbitrary task sequence $\T = (T_1, \dotsc, T_n)$.  
Partition this into two sets of tasks: $\T_{\manualLetter}$ containing all (singleton) tasks that are completed manually, and $\T_{\AIletter}$ containing all other tasks (consisting of AI-augmented singletons and AI chains).  
We first claim that 
\[ FI \leq \sum_{T_j \in \T_{\AIletter}}\left(1 + \sum_{s_i \in T_j}(1-\alpha^{d_i})(1+\minTime{i})\right) + \sum_{(s_i) \in \T_{\manualLetter}}\manualTime{i}. \]

To see why, first note that for any human-executed task $(s_i) \in \T_{\manualLetter}$, either the task fails in the realization of $F$ (in which case it contributes $\minTime{i} \leq \manualTime{i}$ to $FI$), or it succeeds (in which case it appears in some connected component $C_j$, contributing at most $\manualTime{i}$ to the component's weight $\omega(C_j)$).

Next consider the AI-assisted tasks.  
For any $T_j \in \T_{\AIletter}$, consider the set of steps in $T_j$ that fail in any given realization of $F$. 
If no steps in $T_j$ fail, then $T_j$ collectively contributes at most $1$ to the realized fragmentation (since all $s_i \in T_j$ lie in the same connected component of non-failed tasks).  
Otherwise, each failed task $s_i \in T_j$ contributes $\minTime{i}$ to the realized fragmentation (for itself, recalling that $\minTime{i}$ is the minimal cost of executing $s_i$ on its own) plus an additional value of at most $1$ for the weight of a potential new connected component of non-failed tasks.  
As each task $s_i \in T_j$ fails independently with probability $(1-\alpha^{d_i})$, we get that the contribution of tasks in $T_j$ to the fragmentation index is at most $1 + \sum_{s_i \in T_j}(1-\alpha^{d_i})(1+\minTime{i})$ as claimed.

On the other hand, recall that if we take $\T$ to be the cost-minimizing task sequence, then
\[ OPT = \sum_{T_j \in \T_{\AIletter}} \alpha^{-d(T_j)} + \sum_{(s_i) \in \T_{\manualLetter}} \manualTime{i} \] 
by definition, where we write $d(T_j) = \sum_{s_i \in T_j}d_i$ for the total difficulty of steps in task $T_j$.

Comparing our expression for $OPT$ with our bound on $FI$, we see that each task in $\T_{\manualLetter}$ contributes the same amount to each, whereas each $T_j \in \T_{\AIletter}$ contributes $1+\sum_{s_i \in T_j}(1-\alpha^{d_i})(1+\minTime{i})$ to the former and $\alpha^{-d(T_j)}$ to the latter.  
We will show that, for each $T_j \in \T_{\AIletter}$, $1 + \sum_{s_i \in T_j}(1-\alpha^{d_i})(1+\minTime{i}) \leq \frac{5}{4} \alpha^{-d(T_j)}$, which will complete the proof.

First, since $\minTime{i} \leq \alpha^{-d_i}$ for each $s_i \in \mathcal{S}$, we have 
$1 + \sum_{s_i \in T_j}(1-\alpha^{d_i})(1+\minTime{i}) \leq 1 + \sum_{s_i \in T_j}(1-\alpha^{d_i})(1+\alpha^{-d_i})$.

Next note that for any $x,y\in [0,1]$, $(1-xy)(1+\tfrac{1}{xy}) \geq (1-x)(1+\tfrac{1}{x}) + (1-y)(1+\tfrac{1}{y})$.  
(This can be checked by expanding and taking derivatives.)  
Repeatedly applying this fact, we get that $1 + \sum_{s_i \in T_j}(1-\alpha^{d_i})(1+\alpha^{-d_i}) \leq 1 + (1-\alpha^{d(T_j)})(1+\alpha^{-d(T_j)}) = 1 + \alpha^{-d(T_j)} - \alpha^{d(T_j)}$ for any $T_j \in \T_{\AIletter}$.  
The ratio between this expression and $\alpha^{-d(T_j)}$ is maximized at $\alpha^{-d(T_j)} = 1/2$, achieving a value of $5/4$ as claimed.
\end{proof}

Note that this bound of $5/4$ is tight.  
Suppose we have 3 steps in $\mathcal{S}$: the first and last always succeed, and the second succeeds with probability $1/2$ and has a manual cost of $2$.
Then the optimal policy automates all three together for an expected cost of $2$, whereas the fragmentation index is $(1/2)\times 1 + (1/2) \times (1+2+1) = 5/2$.

\medskip

We next argue that $F$ is not much smaller than $OPT$.  
We begin with an analysis under the assumption that $\manualTime{i} \geq 1$ for all $s_i \in \mathcal{S}$.  
That is, for each step $s_i$, completing the step manually is at least as costly as managing an AI tool to complete the step in a single attempt (with guaranteed success).

\begin{proposition}
\label{prop:FI.lower.bound.4}
    Suppose $\manualTime{i} \geq 1$ for all $s_i \in \mathcal{S}$. 
    Then $FI \geq \tfrac{1}{4}OPT$.
\end{proposition}
\begin{proof}

    Consider the following task sequence.  
    Beginning with the first step $s_1$, consider the maximal contiguous sequence of steps $T$ whose success probability $\alpha^{d(T)}$ is greater than or equal to $1/2$.  
    If that set is empty (i.e., the first step has success probability strictly less than $1/2$) then the first step is set to be completed individually, either manually or augmented, whichever is cheaper.  
    In this case, we'll say the resulting singleton task is completed \emph{individually}. 
    Otherwise, the sequence $T = (s_1, \dotsc, s_\ell)$ (which could still be a singleton task) is added to the task sequence as an AI chain; we call this a \emph{non-individual} task.  
    This process is then repeated beginning with the next step (which is $s_2$ in the former case, or $s_{\ell+1}$ in the latter case), until all steps have been added to a task.  
    Call the resulting task sequence $\T$.

    \medskip
    
    Let $ALG$ denote the cost of the task sequence $\T$.  
    Note that we must have $ALG \geq OPT$. 
    We will argue that $FI \geq ALG/4$, which will prove the claim.

    \medskip

    We first define some notation.  
    Write $\T_I$ for the set of individual tasks in $\T$ and $\T_{NI}$ for the set of non-individual tasks.
    Note that $\alpha^{d_i} < 1/2$ for all $(s_i) \in \T_I$, by construction.  
    Note also that $\T_{\manualLetter} \subseteq \T_I$ (recalling that $\T_{\manualLetter}$ is the set of all singleton tasks executed manually), and this containment may be strict if $\T_I$ contains singleton tasks that are augmented.  
    We emphasize that $\T_{NI}$ may still include singleton tasks, but any singleton task $(s_i) \in \T_{NI}$ must have the property that $\alpha^{d_i} \geq 1/2$.
 
    If the final task $s_m$ from $\mathcal{S}$ is contained in some $T_j \in \T_{NI}$, we say that $T_j$ is the \emph{terminal} task.  
    All other tasks in $\T_{NI}$ are non-terminal tasks.  Note that if the final step is a singleton task in $\T_I$ then no task is terminal.  
    For each non-terminal task $T_j \in \T_{NI}$, we'll write $\overline{T}_j$ to denote $T_j$ plus the first step that follows $T_j$. 
    For example, if $T_j = (s_i, \dotsc, s_\ell)$, then $\overline{T}_j = (s_i, \dotsc, s_\ell, s_{\ell+1})$.
    The set $\T_{NI}$ then has the property that for each task $T_j \in \T_{NI}$ we have $\alpha^{d(T_j)} \geq 1/2$, and for each non-terminal $T_j \in \T_{NI}$ we have $\alpha^{(d(\overline{T}_j))} < 1/2$.

    \medskip

    We are now ready to relate $FI$ and $ALG$.  
    Let $k = |\T_{NI}|$ be the number of non-individual tasks. 
    Note first that 
    \[ ALG = \sum_{T_j \in \T_{NI}} \alpha^{-d(T_j)} + \sum_{(s_i) \in \T_I} \minTime{i} \leq 2k + \sum_{(s_i) \in \T_I} \minTime{i}. \]  
    This is because $\alpha^{d(T_j)} \geq 1/2$ for each $T_j \in \T_{NI}$, and hence $\alpha^{-d(T_j)} \leq 2$.

    \medskip

    Next we bound $FI$.  
    Note that the assumption $\manualTime{i} \geq 1$ implies that, in any realization of the connected components of non-failed steps $\mathcal{C} = (C_1, \dotsc, C_\ell)$, $\omega(C_j) = 1$ for all $j$. 
    This implies that the realized fragmentation is equal to $1$ plus, for each $s_i \in F$ (i.e., each step $s_i$ that fails), a contribution of $\minTime{i}$ plus an additional $1$ in the event that $i > 1$ and the task immediately preceding $s_i$ did not fail.  
    (This additional cost of $1$ accounts for creating a new connected component of non-failed steps ending with $s_{i-1}$.)
    
    We claim that $FI \geq k/2 + \frac{1}{2}\sum_{(s_i) \in \T_I} \minTime{i}$.  
    We will show this by charging  the realized fragmentation to tasks in $\T_{NI}$ and $\T_I$, so that each $T_j \in \T_{NI}$ is charged at least $1$ with probability at least $1/2$, and each $(s_i) \in \T_I$ is charged at least $\minTime{i}$ with probability at least $1/2$.
    
    To see this, let $E_j$ denote the event that some step in $\overline{T}_j$ fails, for each non-terminal $T_j \in \T_{NI}$.  
    Note that $E_j$ occurs with probability at least $1/2$, since $\alpha^{d(\overline{T}_j)} < 1/2$. 
    Also, if $E_j$ occurs, then either a step $s_i \in T_j$ fails, or no step in $T_j$ fails but the step immediately following $T_j$ does.  
    In the former case, we will charge the $\minTime{i} \geq 1$ contribution of $s_i$'s failure to $T_j$.  
    In the latter case, it must be that the step immediately preceding the failed task did not fail; so we will charge the additional contribution of $1$ from $s_i$'s failure (as described above) to $T_j$.  
    Additionally, if there is a terminal $T_j$ in $\T_{NI}$, then we charge the baseline $1$ from the calculation of $FI$ to that terminal $T_j$.  
    Aggregating over all these cases, we have that at least $1$ is charged to each set $T_j \in \T_{NI}$ with probability at least $1/2$.  
    Finally, for each $(s_i) \in \T_I$ that fails (which occurs with probability at least $1/2$, by construction), we charge the $\minTime{i}$ portion of its contribution to task $(s_i)$.  
    We conclude that $FI \geq k/2 + \frac{1}{2}\sum_{(s_i) \in \T_I}\minTime{i}$ as claimed.

    But now since $FI \geq k/2 + \frac{1}{2}\sum_{(s_i) \in \T_I}\minTime{i}$ and $ALG \leq 2k + \sum_{(s_i) \in \T_I}\minTime{i}$, we conclude $FI \geq ALG/4$ as claimed.
\end{proof}

Can the approximation factor in Proposition~\ref{prop:FI.lower.bound.4} be improved? 
While we do not have a matching lower bound of $4$, we do know that the approximation factor is at least $\tfrac{2}{3}(2+\sqrt{2}) \approx 2.276$.  
Here is an example. 
Suppose the problem instance is a sequence of $m$ steps, each with success probability $1/\sqrt{2}$ and manual cost $\sqrt{2}$.  
The task sequence described in the proof of Proposition~\ref{prop:FI.lower.bound.4} then handles each task separately, for a total cost of $m \sqrt{2}$.  
The optimal task sequence must therefore perform at least this well.  
The fragmentation index is such that each task contributes $\sqrt{2}$ if it fails, plus $1$ if the preceding task did not fail, for a total contribution of $(1-1/\sqrt{2})(\sqrt{2} + 1(1/\sqrt{2}) = \tfrac{3}{2}(\sqrt{2}-1) \approx 0.6213$.  
As $m$ grows large, the ratio between $m\sqrt{2}$ and $1 + m \times 0.6213$ approaches $\tfrac{2}{3}(2+\sqrt{2}) \approx 2.276$.

\bigskip

We can relax the assumption in Proposition~\ref{prop:FI.lower.bound.4} that $\manualTime{i} \geq 1$ for all $s_i$, but at the cost of a worse approximation factor.

\begin{proposition}
\label{prop:FI.lower.bound.8}
    For general $\manualTime{i}$ (i.e., allowing $\manualTime{i} < 1$ for some steps $s_i$), $FI \geq \frac{1}{8}OPT$.
\end{proposition}
\begin{proof}
    The argument follows the proof of Proposition~\ref{prop:FI.lower.bound.4}, with the following changes.

    First we make a slight change in the definition of the task sequence $\T$ that defines $ALG$.  
    If a constructed task $T_j$ has the property that $\sum_{s_i \in T_j} \manualTime{i} < 1$ (which implies it is cheaper to run all steps of $T_j$ manually than as an AI chain that is guaranteed to succeed), we switch to running the tasks of $T_j$ manually in the task sequence.  
    In our analysis, we still consider $T_j$ to be a set in $\T_{NI}$; but its cost is taken to be $\sum_{s_i \in T_j} \manualTime{i}$.

    This modification changes our upper bound on the total cost of $ALG$ to the following:
    \[ ALG \leq 2 \sum_{T_j \in \T_{NI}} \min\left\{1, \sum_{s_i \in T_j} \manualTime{i} \right\} + \sum_{ (s_i) \in \T_I} \minTime{i}. \]

    Then, to bound $FI$, we claim that 
    \[ FI \geq \frac{1}{4}\sum_{T_j \in \T_{NI}} \min\left\{1, \sum_{s_i \in T_j} \manualTime{i} \right\} + \frac{1}{2}\sum_{(s_i) \in \T_I} \minTime{i} \]
    which would prove the claim.

    To see why this bound on $FI$ holds, we employ a charging argument as before.  
    Recall that for each non-terminal $T_j \in \T_{NI}$, the event $E_j$ occurs with probability at least $1/2$.  
    When it does, we charge to $T_j$ the contribution $\minTime{i}$ to $FI$ from any $s_i \in F$ that lies in $T_j$.  
    Furthermore, when $E_j$ occurs, for each connected component $C_\ell$ that intersects $T_j$ we additionally charge the contribution $\omega(C_\ell)$ from $FI$ to $T_j$.   
    Crucially, the contribution from each $C_\ell$ can be charged to at most two different tasks $T_j$ in this way: once for the leftmost $T_j$ that it intersects, and once for the rightmost $T_j$ that it intersects.  
    (The reason is that if some $T_j$ is a subset of $C_\ell$, then by definition no step of $T_j$ fails and hence $E_j$ did not occur.  
    So this charging can only occur for a task $T_j$ that intersects $C_\ell$ but is not a subset of $C_\ell$, of which there are at most two.)  
    Moreover, this charging to $T_j$ adds up to a  total value of at least $\min\left\{1, \sum_{s_i \in T_j} \manualTime{i} \right\}$.  
    Since the charging occurs with probability $1/2$ for each $T_j \in \T_{NI}$, and we are at most double-counting each $\omega(C_\ell)$, we conclude that the expected sum of all $\omega(C_\ell)$, plus $\minTime{i}$ for all failed tasks $s_i$ that lie in some $T_j \in \T_{NI}$, is at least $1/4$ of $\sum_{T_j \in \T_{NI}}\min\left\{1, \sum_{s_i \in T_j} \manualTime{i} \right\}$.

    The charging for $(s_i) \in \T_I$ remains unchanged: the value $\minTime{i}$ is charged in the event that step $s_i$ fails, which occurs with probability at least $1/2$.

    Taken together, we obtain our desired $8$ approximation.
\end{proof}

We suspect the factor of $8$ in Proposition~\ref{prop:FI.lower.bound.8} is not tight. 
But, as we now show, the factor is strictly greater than the factor $4$ from Proposition~\ref{prop:FI.lower.bound.4}, so the assumption that $\manualTime{i} \geq 1$ for all $s_i$ is necessary for Proposition~\ref{prop:FI.lower.bound.4} to hold.

Consider a step sequence $\mathcal{S}$ that alternates between (a) steps of manual cost $0$ and success probability $1/\sqrt{2} - \epsilon$ where $\epsilon$ is vanishingly small, and (b) steps of manual cost $1$ and success probability $1$.  
Suppose there are $K > 1$ such pairs.  
The task sequence described in the proof of Proposition~\ref{prop:FI.lower.bound.8} above groups the steps into pairs of two-step tasks, for a total cost of $K \sqrt{2}$ (ignoring $\epsilon$).  
The fragmentation index is $1$ plus $1$ for each task that fails, which is $1 + K(1 - 1/\sqrt{2})$ (again ignoring the impact of $\epsilon$).  
As $K$ grows large, the ratio tends to $2(\sqrt{2}+1)$ which is strictly greater than $4$.